% 中文字体与引用类西文字体；四个随附 OTF 不完整时整套回退。
\PassOptionsToPackage{no-math}{fontspec}
\usepackage[UTF8,fontset=none]{ctex}
\defaultfontfeatures{Ligatures=TeX}

% 字体策略：优先使用随项目携带的思源宋体/黑体（inkline/fonts/ 目录，无需系统安装）；
% 如果这些 OTF 被删掉，则回退到 TeX 发行版自带的 Fandol。
% 四个文件必须全部存在，避免字体的 Regular / Bold 指向缺失文件。
% 中文没有西文意义上的 italic：斜体统一映射到 Fandol 楷体，
% 粗斜体保持黑体/宋体粗体，避免出现伪斜体。
\newif\ifink@bundledfonts
\ink@bundledfontstrue
\newcommand{\ink@checkfont}[1]{\IfFileExists{inkline/fonts/#1}{}{\ink@bundledfontsfalse}}
\forcsvlist{\ink@checkfont}{SourceHanSerifCN-Regular.otf,SourceHanSerifCN-Bold.otf,SourceHanSansCN-Regular.otf,SourceHanSansCN-Bold.otf}
\ifink@bundledfonts
  % 完整思源字体：衬线、无衬线与中文等宽族
  \setCJKmainfont{SourceHanSerifCN-Regular.otf}[Path=inkline/fonts/,
    BoldFont=SourceHanSerifCN-Bold.otf,
    ItalicFont=FandolKai-Regular.otf,ItalicFeatures={Path=},
    BoldItalicFont=SourceHanSerifCN-Bold.otf]%
  \setCJKsansfont{SourceHanSansCN-Regular.otf}[Path=inkline/fonts/,
    BoldFont=SourceHanSansCN-Bold.otf,
    ItalicFont=FandolKai-Regular.otf,ItalicFeatures={Path=},
    BoldItalicFont=SourceHanSansCN-Bold.otf]%
  % 等宽场景（代码、路径）用思源黑体，保证中文可读
  \setCJKmonofont{SourceHanSansCN-Regular.otf}[Path=inkline/fonts/,
    BoldFont=SourceHanSansCN-Bold.otf,
    ItalicFont=FandolKai-Regular.otf,ItalicFeatures={Path=},
    BoldItalicFont=SourceHanSansCN-Bold.otf]%
\else
  % 思源字体不完整：整套回退到 Fandol
  \setCJKmainfont{FandolSong-Regular.otf}[BoldFont=FandolSong-Bold.otf,
    ItalicFont=FandolKai-Regular.otf,BoldItalicFont=FandolSong-Bold.otf]%
  \setCJKsansfont{FandolHei-Regular.otf}[BoldFont=FandolHei-Bold.otf,
    ItalicFont=FandolKai-Regular.otf,BoldItalicFont=FandolHei-Bold.otf]%
  \setCJKmonofont{FandolFang-Regular.otf}[BoldFont=FandolHei-Bold.otf,
    ItalicFont=FandolKai-Regular.otf,BoldItalicFont=FandolHei-Bold.otf]%
\fi
% 西文保留 Latin Modern（beamer 默认），无需额外设置。

\IfFileExists{texgyretermes-regular.otf}{%
  \newfontfamily\ink@citefont{texgyretermes-regular.otf}[%
    BoldFont=texgyretermes-bold.otf,
    ItalicFont=texgyretermes-italic.otf,
    BoldItalicFont=texgyretermes-bolditalic.otf]%
}{%
  \newfontfamily\ink@citefont{TeXGyreTermesX-Regular.otf}[%
    BoldFont=TeXGyreTermesX-Bold.otf,
    ItalicFont=TeXGyreTermesX-Italic.otf,
    BoldItalicFont=TeXGyreTermesX-BoldItalic.otf]%
}
